\documentclass[11pt,a4paper]{report}

\usepackage{float} %dzieki temu pakietowi rysunki pojawiaja sie tam gdzie chce po uzyciu opcji H w srodowisku figure
\usepackage[font=small,labelfont=bf]{caption}
\usepackage{natbib}
\usepackage{import}
\usepackage[tight]{subfigure}
\usepackage{tikz}
\usepackage{pgfplots}
\usepackage{polski}
\usepackage[cp1250]{inputenc}
\usepackage{indentfirst}
\usepackage{color}
\usepackage{graphics}
\usepackage{geometry}
% ---- PAGE LAYOUT ----
\geometry{a4paper, top=1.5cm, bottom=1.5cm, left=1.5cm,
right=1.5cm, nohead}
\usepackage{amsmath,amssymb,amsfonts}
\begin{document}
\begin{center}
\large{\textbf{PRACA DOMOWA III}}
\end{center}

\noindent
imi� i nazwisko ........................................\\


\noindent
\textbf{Zadanie 1} Oblicz ca�ki nieoznaczone z funkcji wymiernych:\\ 
\textbf{(a)} $\int \frac{x+7}{x^2-x-6} dx$, 
\textbf{(b)} $\int \frac{-x^2-x-3}{x^3-3 x-2} dx$, 
\textbf{(c)} $\int \frac{4 x-16}{x^3-5 x^2+7 x+13} dx$, 
\textbf{(d)} $\int \frac{2 x^2+3 x-3}{x^3+3 x^2+7 x+5} dx$\\

\vspace{0.25cm}
\noindent
\textbf{odpowiedzi}: \\
\textbf{(a)} $2 \ln (3-x)-\ln (x+2) + C$,
\textbf{(b)} $-\frac{1}{x+1}-\ln (x-2) + C$, 
\textbf{(c)} $\frac{1}{2} \ln \left(x^2-6 x+13\right)-\ln (x+1) + C$,\\  
\textbf{(d)} $\frac{3}{2} \ln \left(x^2+2 x+5\right)-\ln (x+1)-\frac{1}{2} \,\text{arctg}\left(\frac{x+1}{2}\right) + C$.\\

\begin{figure}[H]
	\centering
	\begin{tikzpicture}
		\foreach \x in {0,...,34}
		\draw (0.5*\x,0)--(0.5*\x,21cm);
		\foreach \y in {0,...,42}
		\draw (0,0.5*\y)--(17cm,0.5*\y);
	\end{tikzpicture}
\end{figure}

\begin{figure}[H]
	\centering
	\begin{tikzpicture}
		\foreach \x in {0,...,34}
		\draw (0.5*\x,0)--(0.5*\x,27cm);
		\foreach \y in {0,...,54}
		\draw (0,0.5*\y)--(17cm,0.5*\y);
	\end{tikzpicture}
\end{figure}

\newpage

\noindent
\textbf{Zadanie 2}\\ 
\textbf{(a)} Wyznacz pole zawarte pomi�dzy liniami $y = 3x - x^2$ oraz $y = 2x$.\\ 
\textbf{(b)} Wyznacz pole zawarte pomi�dzy liniami $y = x^2 - 4x + 5$ oraz $y = 5 - x^2$.\\ 
\textbf{(c)} Wyznacz obj�to�� obszaru powsta�ego przez obr�t linii $y = 1-x^2$ wok� osi Ox, gdy $-1\leq x\leq 1 $.\\
\textbf{(d)} Wyznacz obj�to�� obszaru powsta�ego przez obr�t linii $y = \frac{1}{x^2}$ wok� osi Ox, gdy $1\leq x\leq 2$.\\


\noindent
\textbf{odpowiedzi}:\\ 
\textbf{(a)} $S = \frac 16$, \textbf{(b)} $S = \frac 83$, \textbf{(c)} $V = \frac{16}{15}\pi$,  \textbf{(d)} $V = \frac{7}{24}\pi$.\\

\begin{figure}[H]
	\centering
	\begin{tikzpicture}
		\foreach \x in {0,...,34}
		\draw (0.5*\x,0)--(0.5*\x,18cm);
		\foreach \y in {0,...,36}
		\draw (0,0.5*\y)--(17cm,0.5*\y);
	\end{tikzpicture}
\end{figure}

\begin{figure}[H]
	\centering
	\begin{tikzpicture}
		\foreach \x in {0,...,34}
		\draw (0.5*\x,0)--(0.5*\x,27cm);
		\foreach \y in {0,...,54}
		\draw (0,0.5*\y)--(17cm,0.5*\y);
	\end{tikzpicture}
\end{figure}

\newpage

\pagebreak

\noindent
\textbf{Zadanie 3} Wyznacz ekstrema lokalne funkcji:\\
\textbf{(a)} $z = -x^2+4 x y+10 x-8 y^2-4 y+3$, 
\textbf{(b)} $z = -3 x^3-8 x^2-2 x y+9 x+y^2-2 y+6$, 
\textbf{(c)} $z = -4 x^2 y+x^2-4 y^3+3 y$,  
\textbf{(d)} $z = \frac{1}{x^2} + \frac{1}{y^2} + 2 x y$\\

\vspace{0.25cm}

\noindent\textbf{odpowiedzi}:\\ 
\textbf{(a)} punkty krytyczne: $P(9,2)$, $H(P) = 16$, $z''_{xx}(P) = -2$\\ 
\textbf{(b)} punkty krytyczne: $P_1\left(-\frac 73,-\frac 43\right)$, 
$P_2\left(\frac 13,\frac 43\right)$, $H(P_1) = 48$, $z''_{xx}(P_1) = 26$, $H(P_2) = -48$\\ 
\textbf{(c)} punkty krytyczne: $P_1\left(0,-\frac 12\right)$,  
$P_2\left(0,\frac 12\right)$, $P_3\left(-\frac 34,\frac 14\right)$,   
$P_4\left(\frac 34,\frac 14\right)$ , $H(P_1) = 72$, $z''_{xx}(P_1) = 6$, $H(P_2) = 24$, $z''_{xx}(P_2) = -2$, $H(P_3) = -36$, $H(P_4) = -36$\\ 
\textbf{(d)} punkty krytyczne: $P_1(-1,-1)$, 
$P_2(1,1)$, $H(P_1) = 32$, $z''_{xx}(P_1) = 6$, $H(P_2) = 32$, $z''_{xx}(P_2) = 6$\\

\begin{figure}[H]
	\centering
	\begin{tikzpicture}
		\foreach \x in {0,...,34}
		\draw (0.5*\x,0)--(0.5*\x,18cm);
		\foreach \y in {0,...,36}
		\draw (0,0.5*\y)--(17cm,0.5*\y);
	\end{tikzpicture}
\end{figure}

\begin{figure}[H]
	\centering
	\begin{tikzpicture}
		\foreach \x in {0,...,34}
		\draw (0.5*\x,0)--(0.5*\x,27cm);
		\foreach \y in {0,...,54}
		\draw (0,0.5*\y)--(17cm,0.5*\y);
	\end{tikzpicture}
\end{figure}

\end{document}